\section{Cấu hình và ánh xạ sang JAX}

\subsection{Các block config chính}

Tất cả entrypoint quan trọng đều dựa trên YAML của repo. Các block top-level:

\begin{itemize}[leftmargin=1.5em]
  \item \texttt{stage\_1},
  \item \texttt{stage\_2},
  \item \texttt{transport},
  \item \texttt{sampler},
  \item \texttt{guidance},
  \item \texttt{misc},
  \item \texttt{training},
  \item \texttt{eval}.
\end{itemize}

\subsection{Ma trận script và block config}

\begin{center}
\begin{longtable}{>{\raggedright\arraybackslash}p{0.23\linewidth} *{4}{>{\centering\arraybackslash}p{0.11\linewidth}}}
\toprule
Block & \texttt{src\_jax/train.py} & \texttt{src\_jax/sample.py} & \texttt{src\_jax/sample\_ddp.py} & \texttt{src\_jax/stage1\_sample.py} \\
\midrule
\endhead
\texttt{stage\_1} & yes & yes & yes & yes \\
\texttt{stage\_2} & yes & yes & yes & no \\
\texttt{transport} & yes & yes & yes & no \\
\texttt{sampler} & yes & yes & yes & no \\
\texttt{guidance} & yes & yes & yes & no \\
\texttt{misc} & yes & yes & yes & no \\
\texttt{training} & yes & no & no & no \\
\texttt{eval} & yes & no & no & no \\
\bottomrule
\end{longtable}
\end{center}

\subsection{Block \texttt{stage\_1}}

Adapter JAX không tạo schema mới cho Stage 1. Nó đọc trực tiếp:

\begin{itemize}[leftmargin=1.5em]
  \item \texttt{encoder\_config\_path} hoặc \texttt{encoder\_params.dinov2\_path},
  \item \texttt{pretrained\_decoder\_path},
  \item \texttt{normalization\_stat\_path},
  \item \texttt{encoder\_input\_size}.
\end{itemize}

Các giá trị này được chuyển sang encoder RAE của backend NNX để phục vụ
reconstruction và decode latent khi sampling Stage 2.

Nếu cần tính stat riêng cho dataset mới, có thể dùng:

\begin{lstlisting}[style=rae]
python3 src_jax/build_stage1_stats.py \
  --config <stage1_config> \
  --input <train_imagefolder> \
  --output <stage1_stat.pt> \
  --set stage_1.params.normalization_stat_path=<bootstrap_identity_stat.pt>
\end{lstlisting}

Trong đó file bootstrap identity stat chỉ cần chứa \texttt{mean = 0} và
\texttt{var = 1}. Mục tiêu là để pass tính stat đầu tiên không bị chuẩn hóa theo
dataset khác, nhưng vẫn thỏa contract load stat của backend NNX.

\subsection{Block \texttt{stage\_2}}

Một số key quan trọng:

\begin{itemize}[leftmargin=1.5em]
  \item \texttt{target}: chỉ dùng để suy ra backbone NNX tương ứng;
  \item \texttt{ckpt}: có thể là file PyTorch \texttt{.pt} hoặc thư mục Orbax,
  nhưng trên branch này nó phải tương thích với shape của \texttt{SiTDH};
  \item \texttt{input\_size}, \texttt{patch\_size}, \texttt{in\_channels};
  \item \texttt{hidden\_size}, \texttt{depth}, \texttt{num\_heads}: với
  \texttt{SiTDH} đây là các cặp encoder/decoder của backbone DH;
  \item \texttt{class\_dropout\_prob}: tỉ lệ label dropout cho CFG; với
  notebook CelebA một lớp trên đường JAX thì giá trị này được đặt về
  \texttt{0.0};
  \item các toggle như \texttt{use\_rope}, \texttt{use\_rmsnorm},
  \texttt{use\_swiglu}, \texttt{wo\_shift}, \texttt{use\_pos\_embed}.
\end{itemize}

Với config mặc định của branch này, \texttt{target} sẽ là \texttt{SiTDH} và
adapter JAX sẽ ánh xạ nó sang backend \texttt{lightning\_ddt} trong khi vẫn giữ
\texttt{interface\_class=sit}. Các target DDT legacy cũng đi cùng đường
\texttt{lightning\_ddt}.

\subsection{Block \texttt{guidance}}

Hai mode chính:

\begin{itemize}[leftmargin=1.5em]
  \item \texttt{cfg}: adapter dùng chính model làm conditional và unconditional
  pass, với nhãn null mặc định bằng \texttt{num\_classes};
  \item \texttt{autoguidance}: adapter load thêm
  \texttt{guidance.guidance\_model} làm guide network.
\end{itemize}

\subsection{Block \texttt{training}}

Một số key được ánh xạ trực tiếp:

\begin{lstlisting}[style=rae]
training:
  global_seed: 0
  epochs: 1400
  global_batch_size: 1024
  ema_decay: 0.9995
  num_workers: 16
  prefetch_factor: 4
  log_rae_latent_stats: true
  log_activation_stats: true
  log_every: 100
  ckpt_every: 5000
  sample_every: 10000
  base_lr: 0.0002
  final_lr: 0.00002
  beta: [0.9, 0.95]
  wd: 0.0
  schedule_type: linear
\end{lstlisting}

Lưu ý:

\begin{itemize}[leftmargin=1.5em]
  \item nếu config chỉ cho \texttt{epochs}, adapter sẽ suy ra
  \texttt{total\_steps} từ \texttt{num\_train\_samples};
  \item optimizer mặc định vẫn là AdamW;
  \item scheduler được chuẩn hóa về các mode đơn giản như
  \texttt{constant}, \texttt{linear}, \texttt{cosine};
  \item trên đường JAX, \texttt{prefetch\_factor} điều khiển độ sâu prefetch
  theo từng worker cho host-side train loader khi \texttt{num\_workers > 0};
  \item nếu bật \texttt{log\_rae\_latent\_stats=true}, adapter sẽ log
  \texttt{train\_rae\_latent\_*};
  \item nếu bật \texttt{log\_activation\_stats=true}, adapter sẽ log
  RMS/variance cho output Stage 2 và từng block; DH sẽ log theo tên
  \texttt{train\_sitdh\_*}.
\end{itemize}

\subsection{Block \texttt{eval}}

Đường JAX giờ dùng block \texttt{eval} cho cả validation loss lẫn online FID.
Các key \texttt{data\_path}, \texttt{eval\_every}, \texttt{batch\_size},
\texttt{num\_workers}, \texttt{prefetch\_factor}, \texttt{max\_batches}, và
\texttt{eval\_model} sẽ điều khiển held-out validation loop trong
\texttt{src\_jax/train.py}. Nếu không khai báo \texttt{eval.num\_workers} hoặc
\texttt{eval.prefetch\_factor}, adapter JAX sẽ fallback về giá trị tương ứng ở
block \texttt{training}. Với
\texttt{fid\_ref}, nên ưu tiên file \texttt{.pkl}/\texttt{.pickle} được build
bởi \texttt{src\_jax/build\_fid\_stats.py}; nếu đưa file \texttt{.npz} cũ vào,
adapter vẫn tự chuyển nó sang định dạng pickle mà backend NNX hiểu được.
Trên đường JAX, metric \texttt{FID-4K (cfg=...)} mặc định chấm EMA; nếu bật
\texttt{fid\_eval\_model}, adapter sẽ log thêm
\texttt{FID-4K/model (cfg=...)} cho online model mà không thay metric EMA mặc
định.

\subsection{Override từ CLI}

Các entrypoint JAX cho phép override nhanh bằng:

\begin{lstlisting}[style=rae]
python3 src_jax/train.py \
  --config <config> \
  --set training.global_batch_size=256 \
  --set training.prefetch_factor=4 \
  --set eval.prefetch_factor=4 \
  --set training.log_rae_latent_stats=true \
  --set training.log_activation_stats=true \
  --set guidance.scale=1.5
\end{lstlisting}

Mục tiêu là giữ đúng một schema YAML, nhưng vẫn cho phép chỉnh nhanh như CLI.
